\chapter*{INTRODUCCIÓN}
\addcontentsline{toc}{chapter}{INTRODUCCIÓN}

Para las organizaciones de servicios tecnológicos, ordenar la atención de incidencias y controlar los activos informáticos es clave para sostener la continuidad operativa y cumplir con cada cliente. Cuando el registro de fallas o el seguimiento de equipos depende de medios informales o planillas aisladas, los tiempos de respuesta se vuelven impredecibles y se pierde visibilidad operativa.

En J\&P Periféricos S.A.C., el área de soporte técnico atiende requerimientos y gestiona el inventario de activos TI de empresas clientes externas, sus sucursales y trabajadores. En la operación diaria, los pedidos llegaban por llamadas y WhatsApp sin registro formal; los datos de equipos y licencias estaban en hojas de Excel dispersas y desactualizadas; las visitas \textit{in-situ} se coordinaban de forma empírica con cruces de horario; y faltaba validación automática de planes y Acuerdos de Nivel de Servicio (SLA), así como indicadores formales (KPIs). Esto generaba retrasos, búsquedas manuales, reprocesos, pérdidas operativas estimadas y dificultades para respaldar la calidad del servicio.

Para superar estas limitaciones, este proyecto de innovación desarrolla e implementa un Sistema Web Help Desk Centralizado en J\&P Periféricos S.A.C. enfocado en empresas clientes externas. Comprende la recepción y asignación centralizada de tickets, el expediente e inventario multinivel (empresa, sucursal, área y trabajador), la programación de visitas \textit{in-situ}, la validación digital de planes y SLA, el escalamiento estandarizado y los tableros de KPIs (MTTA, MTTR, FCR, CSAT), limitándose al soporte técnico y control de activos TI sin incluir cotizaciones, ventas de hardware o software, facturación ni pasarelas de pago.

El documento consta de ocho capítulos: Capítulo I (generalidades de la empresa); Capítulo II (problema, objetivos, antecedentes y marco teórico); Capítulo III (situación actual, efecto del problema, Ishikawa y Pareto); Capítulo IV (metodología, plan por \textit{sprints} y evidencias de ejecución); Capítulo V (costos de implementación); Capítulo VI (beneficio/costo y rentabilidad); Capítulo VII (conclusiones); y Capítulo VIII (recomendaciones), además de referencias bibliográficas y anexos.